\documentclass[11pt,a4paper]{article}

\usepackage[utf8]{inputenc}
\usepackage[T1]{fontenc}
\usepackage[french]{babel}
\usepackage{lmodern}
\usepackage{amsmath,amssymb}
\usepackage{graphicx}
\usepackage{booktabs}
\usepackage{array}
\usepackage{xcolor}
\usepackage[margin=2.4cm]{geometry}
\usepackage{enumitem}
\usepackage{hyperref}
\usepackage{caption}
\hypersetup{colorlinks=true, linkcolor=blue!55!black, citecolor=blue!55!black,
            urlcolor=blue!55!black}

\definecolor{good}{RGB}{20,120,40}
\definecolor{bad}{RGB}{170,30,30}
\newcommand{\ok}[1]{\textcolor{good}{#1}}
\newcommand{\ko}[1]{\textcolor{bad}{#1}}

\title{\textbf{Estimation de la fréquence cardiaque par rPPG\\
robuste aux carnations foncées}\\[4pt]
\large Comparaison de méthodes classiques et profondes, fine-tuning de PhysNet,\\
approche par régions + CNN~1D, et pistes d'amélioration}
\author{Projet \texttt{mp-rppg-eval}}
\date{\today}

\begin{document}
\maketitle

\begin{abstract}
\noindent
Ce rapport synthétise nos travaux sur l'estimation sans contact de la fréquence
cardiaque (FC) à partir d'une vidéo du visage (\emph{remote photoplethysmography},
rPPG), avec un objectif central : la robustesse sur les \textbf{carnations foncées}
(Fitzpatrick~4--6), là où la littérature et les solutions commerciales restent
fragiles. Nous comparons deux familles d'approches --- les méthodes de
\emph{démixage chromatique} (CHROM, POS) et les réseaux profonds spatio-temporels
(PhysNet) --- et nous introduisons deux approches conçues pour la robustesse
inter-carnation : (i)~un \textbf{CNN~1D} appliqué à des signaux multi-régions /
multi-espaces colorimétriques extraits par segmentation faciale profonde (BiSeNet),
et (ii)~\textbf{CHROM-ITA}, une variante de CHROM dont les coefficients de
projection sont conditionnés par l'angle typologique individuel (ITA) de la peau.
Nous décrivons la motivation théorique de chaque choix, le protocole expérimental
(évaluation au \emph{niveau scénario} et politique de confiance fondée sur le
rapport signal/bruit, SNR), les résultats obtenus sur le jeu VitalVideos
(\texttt{DataVital}), et un ensemble de pistes à explorer --- notamment la
\textbf{fusion par sélection SNR}, l'ajout d'un canal \textbf{ballistocardiographique}
(BCG) indépendant de la couleur de peau, et une architecture de décision multi-signal
inspirée de Binah.AI.
\end{abstract}

\tableofcontents
\bigskip

%======================================================================
\section{Contexte et objectif}
%======================================================================

La photopléthysmographie à distance (rPPG) exploite les minuscules variations de
couleur de la peau induites, à chaque battement cardiaque, par la variation du
volume sanguin dans le lit capillaire sous-cutané~\cite{verkruysse2008,
poh2010}. Une caméra RGB ordinaire suffit en principe à reconstruire un signal
pulsatile et à en déduire la FC.

\paragraph{Le problème des carnations foncées.}
L'amplitude du signal rPPG est proportionnelle à la lumière \emph{rétrodiffusée}
après avoir traversé l'épiderme. Or la mélanine, concentrée dans l'épiderme,
absorbe fortement dans le visible (surtout aux courtes longueurs d'onde). Sur une
peau Fitzpatrick~5--6, le rapport signal/bruit (SNR) du signal pulsatile chute
donc mécaniquement, et toutes les méthodes optiques se dégradent~\cite{nowara2020,
dasari2021}. C'est un enjeu d'équité : plusieurs solutions grand public
sous-performent sur les peaux foncées. Notre objectif est de \textbf{quantifier}
cette dégradation sur nos méthodes et de la \textbf{réduire}.

\paragraph{Jeu de données.}
Nous travaillons sur des vidéos au format \emph{VitalVideos}~\cite{toye2023}
(dossier \texttt{DataVital}) : chaque sujet est filmé en plusieurs
\emph{scénarios} ($\sim$20\,s, 30--60\,fps), avec une vérité-terrain PPG
synchronisée (oxymètre de doigt) et des métadonnées (âge, sexe, \textbf{phototype
Fitzpatrick}, environnement). VitalVideos est précisément conçu pour couvrir une
diversité de carnations, ce qui en fait un banc d'essai pertinent. À ce jour
$\sim$89~sujets sont prétraités et exploités pour l'entraînement~; un lot
supplémentaire de 22~sujets jamais vus sert de test \emph{held-out}.

%======================================================================
\section{Méthodes comparées}
%======================================================================

\subsection{Démixage chromatique : CHROM et POS}

Ces méthodes \emph{non supervisées} partent d'un modèle de réflectance de la peau.
On extrait la couleur moyenne RGB d'une région de peau au cours du temps, on la
normalise, puis on projette les trois canaux sur une direction où la pulsation est
maximisée et l'artéfact de mouvement (variation d'intensité spéculaire) annulé.

\begin{itemize}[leftmargin=1.4em]
\item \textbf{CHROM}~\cite{dehaan2013} construit deux signaux chrominance
$X = 3R_n - 2G_n$ et $Y = 1.5R_n + G_n - 1.5B_n$ (coefficients issus d'une
hypothèse de teinte de peau \emph{standardisée}), puis combine
$S = X - \alpha Y$ avec $\alpha$ égalisant les écarts-types. C'est la méthode de
référence, robuste et rapide.
\item \textbf{POS} (\emph{Plane-Orthogonal-to-Skin})~\cite{wang2017} définit un
plan orthogonal à la direction moyenne de la peau dans l'espace RGB temporel et y
projette le signal. Théoriquement mieux fondée que CHROM sur le plan du modèle de
réflectance, souvent plus robuste au mouvement.
\end{itemize}

\noindent\textbf{Limite pour notre problème.} Les coefficients de CHROM/POS sont
\emph{figés} et calibrés implicitement sur une teinte de peau « moyenne ». Sur une
carnation très foncée (faible signal, spectre de réflectance déplacé), cette
projection fixe n'est plus optimale --- c'est la motivation directe de
CHROM-ITA (§\ref{sec:chromita}).

\subsection{Réseau profond spatio-temporel : PhysNet}

\textbf{PhysNet}~\cite{yu2019} est un encodeur-décodeur \emph{3D-CNN} : il prend
en entrée un cube spatio-temporel (typiquement $128$~images $\times 72\times72$
pixels du visage) et régresse directement la forme d'onde rPPG, image par image.
Les convolutions 3D capturent conjointement la texture spatiale et la dynamique
temporelle, ce qui permet en principe d'apprendre des indices que CHROM/POS
ignorent (micro-textures, distribution spatiale de la pulsation).

\paragraph{Prétraitement \emph{DiffNormalized}.} Comme dans la lignée
DeepPhys/MTTS-CAN~\cite{chen2018,liu2020}, l'entrée est la différence normalisée
entre images consécutives, qui accentue les variations pulsatiles et atténue
l'éclairage statique. \emph{Revers de la médaille}~: sur une vidéo smartphone
fortement compressée (H.264), cette différence amplifie les artéfacts de
compression temporelle --- une limite que nous avons observée et documentée.

\paragraph{Pourquoi fine-tuner.} Les poids publics de PhysNet sont entraînés sur
des jeux à dominante de peaux claires (SCAMPS synthétique, UBFC, PURE). Un
fine-tuning sur des données diversifiées (VitalVideos) vise à recalibrer le réseau
pour les carnations sous-représentées.

\subsection{Approche par régions + CNN 1D}
\label{sec:cnn1d}

Idée (proche de la famille \emph{MSTmap}/RhythmNet~\cite{niu2019,niu2020}) :
plutôt qu'un 3D-CNN coûteux sur des pixels bruts, on \textbf{réduit} d'abord la
vidéo à un petit ensemble de \emph{signaux 1D} robustes, puis on les traite par un
réseau léger.

\begin{enumerate}[leftmargin=1.4em]
\item \textbf{Segmentation faciale profonde} par \textbf{BiSeNet}~\cite{yu2018}
(\emph{face parsing}) : on isole précisément la peau et on découpe le visage en
\textbf{23~régions} (7~régions anatomiques --- front, joues, nez, menton... ---
+ une grille $4\times4$).
\item Pour chaque région et chaque image, on calcule la couleur moyenne dans
\textbf{9~canaux} couvrant 3~espaces colorimétriques (RGB, YUV, Lab), soit
$23\times9 = 207$~signaux temporels.
\item Un \textbf{CNN~1D} dilaté ($\sim$113\,k~paramètres, dilatations $1,2,4,8$)
régresse la forme d'onde rPPG à partir de ces 207~canaux.
\end{enumerate}

\noindent\textbf{Justification.} (i)~La segmentation profonde garantit qu'on
n'agrège que des pixels de peau (vs une boîte rectangulaire qui inclut cheveux,
fond, vêtements)~; (ii)~la \emph{redondance} multi-régions / multi-espaces donne
au réseau des combinaisons linéaires riches pour isoler la pulsation du bruit,
\emph{y compris} des combinaisons mieux adaptées à une carnation donnée
--- ce qu'un CHROM à coefficients fixes ne peut pas faire~; (iii)~le modèle est
minuscule et rapide, donc déployable. \textbf{Risque identifié}~: avec 207~canaux
et peu de sujets, le sur-apprentissage guette (cf.~§\ref{sec:results}).

\subsection{CHROM-ITA : démixage conditionné par la carnation}
\label{sec:chromita}

Pour lever la rigidité de CHROM tout en gardant son interprétabilité, nous
conditionnons ses coefficients sur la \textbf{carnation mesurée}. L'\textbf{ITA}
(\emph{Individual Typology Angle})~\cite{chardon1991} est un descripteur continu
de teinte de peau dérivé de l'espace L*a*b* :
\[
\mathrm{ITA} = \arctan\!\Big(\frac{L^* - 50}{b^*}\Big)\cdot\frac{180}{\pi},
\]
positif pour les peaux claires, négatif pour les peaux foncées. Un petit MLP
($1\!\to\!8\!\to\!5$) prend l'ITA de la région et produit les coefficients de
projection chromatique, \textbf{initialisés sur les valeurs de De Haan} (de sorte
que le modèle non entraîné $\equiv$ CHROM standard). On apprend ainsi une
\emph{famille continue} de démixages indexée par la carnation.

\subsection{Repères industriels : Binah.AI}
\label{sec:binah}

Binah.AI revendique une mesure multi-constantes (FC, VFC, SpO$_2$, respiration,
stress) robuste à toutes les carnations, depuis une simple caméra. D'après leurs
communications, leur pipeline repose sur~: (i)~le \textbf{TOI} (\emph{Transdermal
Optical Imaging}), c.-à-d. la même famille « couleur » que CHROM/POS~; (ii)~un
module appris de \textbf{qualité de signal (SQA)} qui \textbf{rejette} les fenêtres
trop bruitées et ne renvoie une mesure que lorsqu'elle est jugée fiable (d'où le
fameux « measuring\dots » qui se prolonge en conditions difficiles)~; (iii)~de
\emph{très gros} jeux d'entraînement multi-ethniques pour calibrer la conversion
signal~$\rightarrow$~constante vitale. \textbf{Enseignement pour nous}~: leur
avantage n'est pas un algorithme unique « anti-mélanine », mais une
\textbf{architecture de décision} (multi-signal + rejet par qualité). Nos briques
SNR (§\ref{sec:snr}) vont exactement dans ce sens.

%======================================================================
\section{Protocole expérimental}
%======================================================================

\subsection{Audit du prétraitement (avant ré-entraînement)}

Avant tout entraînement « propre » à grande échelle, nous avons audité la chaîne
et corrigé trois défauts qui faussaient les résultats~:
\begin{itemize}[leftmargin=1.4em]
\item \textbf{Augmentation temporelle de PhysNet} : le retournement temporel de
l'entrée n'inversait pas la cible $y$ $\Rightarrow$ paires (entrée, label)
désalignées (corruption de label). \emph{Corrigé} : $x$ et $y$ retournés
conjointement.
\item \textbf{Double filtrage passe-bande} dans le pipeline CNN~1D (label déjà
filtré à l'extraction, puis re-filtré à l'entraînement). \emph{Corrigé}.
\item \textbf{Fréquences d'images hétérogènes} (40/60\,fps à l'entraînement,
$\sim$30\,fps au déploiement). \emph{Corrigé} : ré-échantillonnage de toutes les
séquences à \textbf{30\,fps} (image la plus proche) pour homogénéiser la
dynamique temporelle apprise.
\end{itemize}
La détection/suivi du visage a par ailleurs été basculée de Haar Cascade vers
\textbf{MediaPipe FaceMesh} en mode \emph{tracking} (landmarks propagés d'image en
image), plus stable et évitant qu'une main (geste « facepalm ») soit prise pour le
visage --- les régions de peau étant de toute façon resegmentées par BiSeNet.

\subsection{Évaluation au niveau scénario (et non par fenêtre)}

Une erreur méthodologique fréquente consiste à mesurer la FC sur des fenêtres
courtes (128~images $\approx$~2\,s) et à moyenner. Sur 2\,s, la résolution
fréquentielle de la FFT est grossière et la FC estimée est instable. Nous
reconstruisons donc le signal \textbf{complet du scénario} ($\sim$20\,s) puis
estimons \textbf{une} FC par FFT. Cela change radicalement les chiffres
(p.~ex. une MAE PhysNet passant de $\sim$19,8 par fenêtre à $\sim$7--8 au niveau
scénario) et constitue la seule comparaison équitable entre méthodes.

\subsection{Politique de confiance fondée sur le SNR}
\label{sec:snr}

Plutôt que de toujours renvoyer une FC, nous calculons pour chaque estimation un
\textbf{SNR aveugle} : énergie spectrale dans une bande étroite autour du pic FC
estimé, rapportée à l'énergie hors-bande dans $[0{,}7,\,2{,}5]$\,Hz. Nous
appliquons alors la politique (inspirée du SQA de Binah, §\ref{sec:binah})~:
\[
\boxed{~\text{SNR} > -1~\Rightarrow~\ok{VALIDÉ}\quad
-3 \le \text{SNR} \le -1~\Rightarrow~\text{douteux}\quad
\text{SNR} < -3~\Rightarrow~\ko{REJETÉ}~}
\]
Empiriquement, les échecs se concentrent vers $\sim\!-3$\,dB et les succès vers
$\sim\!+1{,}5$\,dB : le SNR est donc un \emph{signal de confiance} exploitable.
\emph{Réserve importante} : il n'est pas infaillible (cf. CHROM-ITA, §\ref{sec:results}).

%======================================================================
\section{Résultats}
\label{sec:results}
%======================================================================

\subsection{Choix de la base PhysNet à fine-tuner}

Nous avons évalué les \textbf{5~jeux de poids publics} (SCAMPS, UBFC, PURE,
MA-UBFC, BP4D) sur nos données. \textbf{PURE} généralise le mieux et a été retenu
comme base. (SCAMPS, purement synthétique, marche bien sur UBFC mais mal sur nos
vidéos~; UBFC seul généralise encore moins.)

\subsection{Stratégie de fine-tuning : geler ou non ?}

Nous avons comparé quatre stratégies de gel sur la base PURE~:
\begin{center}
\begin{tabular}{@{}lll@{}}
\toprule
Variante & Couches entraînées & Résultat \\
\midrule
A & \textbf{tout} le réseau (full fine-tune) & \ok{meilleure} \\
B & gel ConvBlock4--9 + upsampling & moins bon \\
C & B + recalibration BatchNorm & moins bon \\
D & gel du seul goulot (ConvBlock8--9) & moins bon \\
\bottomrule
\end{tabular}
\end{center}
\textbf{Conclusion}~: le \textbf{full fine-tune (A)} domine. Le gel partiel crée
un \emph{goulot de capacité} --- le réseau n'a pas assez de degrés de liberté pour
se recalibrer sur la nouvelle distribution de carnations. On en retient aussi
qu'il est normal qu'une perte par \emph{corrélation de Pearson} (négative) ne
« tende pas vers~0 »~: elle est bornée par le bruit intrinsèque de la
vérité-terrain~; une bonne corrélation peut coexister avec une MAE encore élevée,
les deux mesurant des choses différentes.

\subsection{Comparaison des méthodes (niveau scénario, base propre)}

Sur le découpage de test (graine~42, 89~sujets, 20~scénarios de test), au niveau
scénario, avec la politique SNR~:

\begin{center}
\captionof{table}{Comparaison globale. « Validés » = scénarios avec SNR$>-1$~;
« MAE val. » = MAE sur ces seuls scénarios.}
\begin{tabular}{@{}lccccc@{}}
\toprule
Méthode & MAE (bpm) & médiane & \%\,$<5$\,bpm & Validés (SNR$>\!-1$) & MAE val. \\
\midrule
\textbf{CNN~1D}    & \textbf{6,77} & 0,0 & 70\,\% & --- & --- \\
\textbf{PhysNet}   & 7,03 & 2,64 & 65\,\% & 7/20 & \ok{5,27} \\
\textbf{POS}       & --- & --- & --- & \ok{9/20} & \ok{0,39} \\
\textbf{CHROM-ITA} & 11,43 & --- & --- & --- & \ko{$\sim$26 (cas piégeux)} \\
\bottomrule
\end{tabular}
\end{center}

\noindent\textbf{Lectures clés.}
\begin{itemize}[leftmargin=1.4em]
\item \textbf{Au global}, CNN~1D (6,77) et PhysNet (7,03) sont au coude-à-coude
(chiffres du découpage interne~; les modèles appris ont depuis été ré-entraînés sur
données propres, cf.~§\ref{sec:heldout} pour les scores held-out de référence).
\item \textbf{Sous la politique SNR}, \textbf{POS} domine nettement : 9/20
scénarios validés à \textbf{MAE 0,39}. Il est \emph{conservateur mais fiable}.
\item Le SNR de \textbf{PhysNet est honnête}~: il \emph{rejette} quand il échoue
(p.~ex. erreur de 28\,bpm $\rightarrow$ SNR $-4{,}8$ $\rightarrow$ rejeté). Ses
7~scénarios validés sont propres (MAE 5,27).
\item \textbf{CHROM-ITA est le seul danger}~: il peut être \emph{confiant et faux}
(doublage d'harmonique : il verrouille sur $2\times$\,FC avec un SNR élevé).
Le SNR seul ne suffit pas à le sécuriser $\Rightarrow$ garde anti-harmonique
nécessaire (§\ref{sec:pistes}).
\item Les échecs « peau foncée » sont des \textbf{cas individuels difficiles}
(quelques sujets), pas une dégradation uniforme de toute la classe Fitzpatrick~6.
\end{itemize}

\subsection{Test sur 2 vidéos externes (carnations très foncées)}
\label{sec:extvideo}

Deux vidéos personnelles avec FC de référence mesurée par appareil externe, sur
peaux très foncées (ITA $\le -56$). C'est un test \emph{in vivo} hors du jeu
d'entraînement.

\begin{center}
\captionof{table}{Vidéo A --- réf.~\textbf{54\,bpm}, ITA~$\approx -56$.}
\begin{tabular}{@{}lccc@{}}
\toprule
Méthode & FC estimée & erreur & SNR aveugle \\
\midrule
\textbf{CNN~1D}    & 51,0 & \ok{\textbf{3,0}} & \ok{$-0{,}26$} \\
PhysNet            & 80,9 & 26,9 & $-1{,}28$ \\
CHROM-ITA          & 64,2 & 10,2 & $-3{,}36$ \\
CHROM              & 63,3 & 9,3  & $-3{,}50$ \\
POS                & 96,7 & 42,7 & $-7{,}90$ \\
\midrule
écart-type des FC & \multicolumn{3}{c}{\textbf{15,9\,bpm} (méthodes \emph{divergentes})} \\
\bottomrule
\end{tabular}
\end{center}

\begin{center}
\captionof{table}{Vidéo B --- réf.~\textbf{56\,bpm}, ITA~$\approx -65$.}
\begin{tabular}{@{}lccc@{}}
\toprule
Méthode & FC estimée & erreur & SNR aveugle \\
\midrule
\textbf{CNN~1D}    & 49,2 & 6,8 & \ok{$-0{,}47$} \\
\textbf{POS}       & 56,2 & \ok{\textbf{0,2}} & $-2{,}20$ \\
\textbf{CHROM}     & 56,2 & \ok{\textbf{0,2}} & $-4{,}30$ \\
PhysNet            & 66,1 & 10,1 & $-3{,}64$ \\
CHROM-ITA          & 80,2 & 24,2 & \ko{$-5{,}28$} \\
\midrule
écart-type des FC & \multicolumn{3}{c}{\textbf{10,8\,bpm} (méthodes \emph{regroupées})} \\
\bottomrule
\end{tabular}
\end{center}

\noindent\textbf{Deux régimes opposés --- et c'est l'enseignement central.}
\begin{itemize}[leftmargin=1.4em]
\item \textbf{Vidéo A (divergence).} Les méthodes sont éparpillées~; la moyenne
naïve (75,6\,bpm) est polluée par les fausses. Seule la \emph{sélection par SNR}
(CNN~1D, meilleur SNR) donne la bonne réponse (err~3).
\item \textbf{Vidéo B (accord).} CHROM/POS retrouvent la FC \emph{exacte}
(56,2 vs 56) mais avec un SNR \emph{faible} (signal ténu sur peau Fitz~6)~: le SNR
\emph{sous-note} les bonnes méthodes. Ici la sélection-SNR retiendrait CNN~1D
(err~6,8), tandis que la \textbf{médiane} des méthodes donne 56,2 (err~0,2).
\end{itemize}

\noindent\textbf{Conséquence : une fusion ADAPTATIVE.} Aucune règle fixe (moyenne
seule, ou SNR seul) ne réussit sur les deux vidéos. Nous adoptons donc une fusion
qui \emph{bascule selon l'accord inter-méthodes}~:
\[
\boxed{~\text{écart-type des FC} \le \tau~\Rightarrow~\textbf{consensus} \text{ (médiane)}
\qquad
\text{écart-type} > \tau~\Rightarrow~\textbf{sélection} \text{ (meilleur SNR validé)}~}
\]
avec $\tau \approx 12$\,bpm et une \emph{garde anti-harmonique} (on écarte un
candidat dont la FC vaut $\approx 2\times$ la médiane des autres). Résultats de
cette fusion adaptative~:

\begin{center}
\begin{tabular}{@{}lcccc@{}}
\toprule
Vidéo & réf. & mode déclenché & FC fusionnée & erreur \\
\midrule
A & 54 & sélection ($\rightarrow$ CNN~1D) & 51,0 & \ok{\textbf{3,0}} \\
B & 56 & consensus ($\rightarrow$ médiane) & 56,2 & \ok{\textbf{0,2}} \\
\bottomrule
\end{tabular}
\end{center}

\noindent Cette stratégie est désormais la \textbf{méthode de fusion par défaut}
du pipeline (\texttt{mp\_rppg/fusion.py}, utilisée par \texttt{run\_on\_video.py}).

\subsection{Bilan terrain : 5 vidéos externes, tout le spectre de carnations}
\label{sec:bilanterrain}

En agrégeant toutes les vidéos personnelles testées avec le pipeline corrigé
(modèles ré-entraînés sur données propres, fusion adaptative, CNN~1D ré-inclus),
on couvre un large spectre d'ITA, du clair ($+3$) au très foncé ($-65$). Une de ces
vidéos (\texttt{IMG\_9008}) a été traitée en \textbf{test à l'aveugle} (réf.
révélée \emph{après} la prédiction).

\begin{center}
\captionof{table}{Erreur de la \textbf{fusion adaptative} sur les vidéos externes.
« mode » = branche déclenchée. ($\dagger$ = test à l'aveugle.)}
\begin{tabular}{@{}lcccccl@{}}
\toprule
Vidéo & réf. (bpm) & ITA & carnation & fusion & erreur & mode \\
\midrule
Issa63               & 63 & $+3$  & claire       & 64,1 & \ok{1,1} & consensus \\
IMG\_9008$^\dagger$  & 67 & $-16$ & médiane      & 69,4 & \ok{2,4} & consensus \\
BPM54                & 54 & $-56$ & foncée       & 51,0 & \ok{3,0} & sélection \\
BPM56                & 56 & $-65$ & très foncée  & 56,2 & \ok{0,2} & consensus \\
VideoTest4           & 58 & $-62$ & très foncée  & 62,4 & 4,4 & consensus \\
\bottomrule
\end{tabular}
\end{center}

\noindent\textbf{Résultat} : sur les 5~vidéos, la fusion adaptative reste
\textbf{sous 5\,bpm d'erreur} de l'ITA~$+3$ à l'ITA~$-65$ (4~cas sous 3\,bpm, un à
4,4), en adaptant son mode à chaque vidéo. Le test à l'aveugle confirme l'absence de
réglage \emph{a posteriori}. \emph{Deux réserves honnêtes} :
\begin{itemize}[leftmargin=1.4em]
\item Sur les carnations très foncées, les SNR sont souvent négatifs (signal ténu)
--- la justesse vient alors de l'\textbf{accord inter-méthodes} plus que du SNR, ce
qui justifie d'utiliser \emph{deux} indices de confiance (SNR \emph{et} écart-type).
\item Le cas le moins bon (VideoTest4, erreur 4,4) illustre une faiblesse de la
médiane brute~: deux méthodes (CHROM-ITA, POS) surestimaient \emph{dans le même
sens} (71 et 75\,bpm), tirant la médiane à 62,4 alors que le vrai cluster bas
--- CNN~1D (57,1~; erreur 0,9) et PhysNet (60,6) --- était correct. Une médiane
\textbf{pondérée par le SNR} ou un \textbf{rejet d'outlier} avant consensus
corrigerait ce biais (piste §\ref{sec:pistes}).
\end{itemize}

\paragraph{Garde-fou vidéos 4K.} Les vidéos haute résolution (p.~ex. 2160$\times$3840)
saturent la mémoire si toutes les images sont chargées en pleine résolution. Le
chargeur redimensionne désormais les images à \texttt{max\_dim=720} dès la lecture
(option \texttt{-{}-max-dim} de \texttt{run\_on\_video.py}) --- sans impact sur le
signal rPPG, qui est temporel et non lié à la résolution spatiale.

\subsection{Test \emph{held-out} sur 22 sujets jamais vus (majorité Fitz 5--6)}
\label{sec:heldout}

Test le plus exigeant~: 22~nouveaux sujets VitalVideos (54~scénarios),
\textbf{jamais vus à l'entraînement}, à dominante de carnations foncées. Évaluation
au niveau scénario.

\begin{center}
\captionof{table}{Held-out, niveau scénario, \textbf{après correction du bug de
modèles périmés} (tous les modèles appris ré-entraînés sur l'extraction propre).
« Validés » = SNR$>-1$. La colonne « périmé » rappelle le score erroné avant
correction.}
\begin{tabular}{@{}lccccc|c@{}}
\toprule
Méthode & MAE & médiane & \%\,$<5$ & Validés & MAE val. & (périmé) \\
\midrule
\textbf{PhysNet}   & \ok{\textbf{3,89}} & 0,00 & \ok{\textbf{83\,\%}} & 35/52 & \ok{\textbf{0,25}} & 3,89 \\
\textbf{CNN~1D}    & \ok{\textbf{4,69}} & 0,00 & 80\,\% & 40/54 & 1,27 & \ko{26,60} \\
CHROM-ITA          & 8,27 & 0,00 & 67\,\% & 21/54 & 4,27 & 11,82 \\
CHROM              & 8,69 & 0,88 & 59\,\% & 17/54 & 3,52 & --- \\
POS                & 9,90 & 0,88 & 65\,\% & 23/54 & 3,90 & --- \\
\midrule
Fusion adaptative  & 4,69 & 0,00 & 81\,\% & \multicolumn{2}{c}{consensus×43, sél.×11} & \ko{15,30} \\
\bottomrule
\end{tabular}
\end{center}

\paragraph{(A) Un bug de \emph{modèle périmé}, pas un défaut d'architecture
--- correction d'une conclusion intermédiaire.} Une première lecture avait conclu
que le CNN~1D « collapsait » (MAE 26,6, prédictions quasi-constantes, corrélation
$-0{,}21$ avec la vraie FC) et qu'il fallait réduire sa capacité. \textbf{Le
diagnostic réel est autre.} Les fichiers de poids du CNN~1D et de CHROM-ITA avaient
été entraînés \emph{le matin}, sur une extraction de \texttt{region\_signals}
\textbf{antérieure aux corrections de prétraitement} (double passe-bande sur les
labels + fps mixtes), puis jamais ré-entraînés après la ré-extraction propre de
midi. On évaluait donc des modèles appris sur des labels buggés. \textbf{Preuves}~:
(i)~les dates des poids précèdent la ré-extraction~; (ii)~ré-entraîner la
\emph{même} architecture sur les données propres redonne un modèle sain
(CNN~1D : MAE \textbf{4,69}, corr 0,65~; CHROM-ITA : 11,82~$\rightarrow$~8,27)~;
(iii)~une étude d'ablation des canaux (207 / 69 / 21) montre que \textbf{réduire la
capacité n'aide pas} --- la version 207~canaux ré-entraînée propre est la meilleure
--- donc le problème n'était pas le sur-apprentissage. \emph{Leçon de processus}~:
toujours ré-entraîner les modèles \emph{après} une correction du prétraitement, et
vérifier les horodatages poids~vs~données.

\paragraph{(B) PhysNet domine \emph{en distribution}, le CNN~1D suit de près.} Sur
sujets VitalVideos non vus, PhysNet généralise remarquablement (MAE \textbf{3,89}~;
83\,\% $<5$\,bpm~; 0,25 sur les validés)~; le CNN~1D ré-entraîné est juste derrière
(4,69~; 80\,\%). La \textbf{fusion adaptative} (4,69, désormais avec CNN~1D ré-inclus)
égale le CNN~1D mais ne dépasse pas PhysNet seul~: en distribution, ajouter de
l'optique plus bruitée dilue la bonne estimation profonde.

\paragraph{(C) Réconciliation avec les vidéos externes.} Contradiction apparente~:
sur les 2~vidéos smartphone (§\ref{sec:extvideo}) PhysNet était \emph{mauvais},
ici il est \emph{excellent}. L'explication est le \textbf{décalage de domaine}~: le
held-out est au format VitalVideos (même distribution que l'entraînement~: bon
capteur, éclairage maîtrisé), tandis que les vidéos perso sont du smartphone H.264
(compression temporelle), distribution sur laquelle le prétraitement DiffNormalized
de PhysNet se dégrade. \emph{Leçon opérationnelle}~:
\begin{itemize}[leftmargin=1.4em]
\item \textbf{En distribution} (vidéo type VitalVideos)~: \textbf{faire confiance à
PhysNet} (3,89\,bpm), la fusion n'aide pas.
\item \textbf{Hors distribution} (smartphone/compressé)~: PhysNet n'est plus
fiable~; les méthodes \emph{sans modèle} (CHROM/POS) et la \textbf{fusion
adaptative} deviennent le filet de sécurité.
\end{itemize}
La fusion adaptative garde donc toute sa valeur comme \textbf{robustesse au domaine
inconnu}, mais ne doit pas masquer que, sur sa distribution, PhysNet est le
meilleur estimateur isolé.

%======================================================================
\section{Discussion}
%======================================================================

\paragraph{Aucune méthode unique ne suffit.} Chaque méthode a un régime de
défaillance \emph{différent}~: POS est très précis quand il « voit » le signal
mais se tait souvent~; CNN~1D est le plus stable y compris sur peau foncée mais
peut sur-apprendre~; PhysNet est honnête (bon SNR) mais conservateur~; CHROM-ITA
est risqué. Cette \textbf{complémentarité} est précisément ce qui rend une fusion
intéressante.

\paragraph{Le SNR est un bon proxy de confiance, mais pas une garantie.} Il
sépare bien succès et échecs \emph{en moyenne}, mais le doublage d'harmonique
(CHROM-ITA) montre qu'un signal peut être spectralement « propre » tout en
pointant la mauvaise fréquence. D'où la nécessité de garde-fous physiologiques.

\paragraph{Honnêteté statistique.} Les comparaisons entre sous-groupes
Fitzpatrick doivent se faire à découpage de test \emph{identique} : une « MAE
Fitz~6 = 15\,bpm » qui repose en réalité sur un unique sujet difficile n'est pas
une preuve de biais systématique. Nous avons aussi vérifié l'\textbf{absence de
fuite} entre train et test (GUID uniques) lorsqu'un résultat paraissait « trop
bon ».

%======================================================================
\section{Pistes à explorer}
\label{sec:pistes}
%======================================================================

\begin{enumerate}[leftmargin=1.5em]

\item \textbf{Fusion adaptative (ADOPTÉE, cf.~§\ref{sec:extvideo}).} Désormais
méthode par défaut~: \emph{médiane} quand les méthodes s'accordent (faible
écart-type), \emph{sélection par SNR} quand elles divergent, avec garde
anti-harmonique et sortie possible « mesure non fiable » si tout échoue (logique
Binah). \emph{À faire}~: (a)~calibrer finement le seuil $\tau$ et le SNR de
validation sur le test \emph{held-out}~; (b)~remplacer la médiane brute du mode
consensus par une \textbf{médiane pondérée par le SNR} (ou un \textbf{rejet
d'outlier} préalable), pour éviter qu'une paire de méthodes biaisées dans le même
sens tire le consensus (cas VideoTest4, §\ref{sec:bilanterrain}).

\item \textbf{Garde anti-harmonique.} Avant de valider une FC, vérifier qu'il
n'existe pas un pic d'énergie comparable à FC/2 (sous-harmonique) dans la bande
physiologique~; sinon corriger ou déclasser. Sécurise spécifiquement CHROM-ITA.

\item \textbf{Canal ballistocardiographique (BCG), indépendant de la couleur de
peau.} Le BCG~\cite{balakrishnan2013} mesure les \emph{micro-mouvements} de la
tête dus à l'éjection systolique (recul newtonien de l'afflux carotidien). Étant
\emph{mécanique}, il est \textbf{insensible à la mélanine} --- exactement
complémentaire de l'optique~:
\begin{center}
\begin{tabular}{@{}lcc@{}}
\toprule
 & rPPG optique & BCG (mouvement) \\
\midrule
Source physique        & couleur de peau & déplacement de la tête \\
Sensible à la mélanine & \ko{oui} & \ok{non} \\
Sensible au mouvement volontaire & moyen & \ko{fort} \\
Sensible à l'éclairage & \ko{oui} & \ok{non} \\
\bottomrule
\end{tabular}
\end{center}
\emph{Mise en œuvre} (non supervisée, réutilise nos landmarks)~: suivre le
déplacement vertical de points rigides (arête du nez, front) $\rightarrow$ ACP des
trajectoires $\rightarrow$ composante la plus périodique dans $[0{,}7,2{,}5]$\,Hz
$\rightarrow$ FFT $\rightarrow$ FC$_\text{BCG}$ + son propre SNR, ajoutée comme
votant supplémentaire à la fusion. \emph{Condition}~: tête immobile et caméra
stable (assis/trépied)~; à proscrire en vidéo smartphone tenue à la main.

\item \textbf{Architecture de décision multi-signal type Binah.} Généraliser la
politique SNR en un module de qualité appris (SQA) opérant par fenêtre, capable de
combiner plusieurs constantes (FC, VFC, respiration) et de rejeter explicitement
les segments douteux.

\item \textbf{Plus de données diversifiées.} La principale limite du CNN~1D
(sur-apprentissage à 207~canaux) et du fine-tuning PhysNet est la taille du jeu.
Poursuivre l'intégration de VitalVideos / d'un sous-ensemble africain (Fitz~5--6)
pour densifier les carnations sous-représentées. Étudier aussi une
\textbf{réduction du nombre de canaux} (moins de régions / d'espaces) pour
réduire le sur-apprentissage du CNN~1D.

\item \textbf{Robustesse à la compression.} PhysNet (DiffNormalized) se dégrade
sur vidéo H.264 fortement compressée~; explorer un prétraitement moins sensible
aux artéfacts temporels, ou un entraînement augmenté par compression simulée.

\end{enumerate}

%======================================================================
\section{Conclusion}
%======================================================================

Le test \emph{held-out} (§\ref{sec:heldout}) impose deux conclusions fortes.
\textbf{(1)~En distribution, les deux modèles appris dominent}~: PhysNet fine-tuné
(MAE \textbf{3,89}\,bpm, 83\,\% $<5$\,bpm) en tête, CNN~1D ré-entraîné juste derrière
(4,69), sur des sujets VitalVideos jamais vus à dominante de peaux foncées --- la
fusion adaptative les égale mais ne dépasse pas PhysNet seul.
\textbf{(2)~Un incident de processus a été identifié et corrigé}~: deux modèles
(CNN~1D, CHROM-ITA) avaient été livrés entraînés sur une extraction \emph{antérieure}
aux corrections de prétraitement, faussant lourdement leur évaluation
(CNN~1D~: 26,6 au lieu de 4,7)~; tous les modèles ont été ré-entraînés sur données
propres. La valeur propre de la \textbf{fusion adaptative} se situe \emph{hors
distribution} (vidéo smartphone compressée, §\ref{sec:extvideo}), où PhysNet décroche
et où les méthodes sans modèle prennent le relais. Stratégie de déploiement
recommandée~: \emph{PhysNet en primaire si le signal est de qualité (SNR élevé),
fusion adaptative (CNN~1D/CHROM/POS $\pm$~futur BCG) en secours sinon}, avec rejet
explicite quand toutes les méthodes échouent. La piste la plus prometteuse pour la
robustesse inter-carnation reste l'ajout d'un canal \textbf{BCG}, physiquement
indépendant de la mélanine, tout en restant \emph{honnête} sur les cas impossibles.

%======================================================================
\begin{thebibliography}{99}
\setlength{\itemsep}{2pt}

\bibitem{verkruysse2008} W.~Verkruysse, L.~O.~Svaasand, J.~S.~Nelson.
\emph{Remote plethysmographic imaging using ambient light}. Optics Express,
16(26):21434--21445, 2008.

\bibitem{poh2010} M.-Z.~Poh, D.~J.~McDuff, R.~W.~Picard.
\emph{Non-contact, automated cardiac pulse measurements using video imaging and
blind source separation}. Optics Express, 18(10):10762--10774, 2010.

\bibitem{dehaan2013} G.~de~Haan, V.~Jeanne.
\emph{Robust pulse rate from chrominance-based rPPG}. IEEE Transactions on
Biomedical Engineering, 60(10):2878--2886, 2013. (CHROM)

\bibitem{wang2017} W.~Wang, A.~C.~den~Brinker, S.~Stuijk, G.~de~Haan.
\emph{Algorithmic principles of remote PPG}. IEEE Transactions on Biomedical
Engineering, 64(7):1479--1491, 2017. (POS)

\bibitem{yu2019} Z.~Yu, X.~Li, G.~Zhao.
\emph{Remote photoplethysmograph signal measurement from facial videos using
spatio-temporal networks}. BMVC, 2019. (PhysNet)

\bibitem{chen2018} W.~Chen, D.~McDuff.
\emph{DeepPhys: Video-based physiological measurement using convolutional
attention networks}. ECCV, 2018.

\bibitem{liu2020} X.~Liu, J.~Fromm, S.~Patel, D.~McDuff.
\emph{Multi-task temporal shift attention networks for on-device contactless
vitals measurement (MTTS-CAN)}. NeurIPS, 2020.

\bibitem{niu2019} X.~Niu, H.~Han, S.~Shan, X.~Chen.
\emph{RhythmNet: End-to-end heart rate estimation from face via spatial-temporal
representation}. IEEE Transactions on Image Processing, 2019.

\bibitem{niu2020} X.~Niu, et al.
\emph{Video-based remote physiological measurement via cross-verified feature
disentangling (MSTmap)}. ECCV, 2020.

\bibitem{yu2018} C.~Yu, J.~Wang, C.~Peng, C.~Gao, G.~Yu, N.~Sang.
\emph{BiSeNet: Bilateral segmentation network for real-time semantic
segmentation}. ECCV, 2018. (face parsing)

\bibitem{balakrishnan2013} G.~Balakrishnan, F.~Durand, J.~Guttag.
\emph{Detecting pulse from head motions in video}. CVPR, 2013. (BCG)

\bibitem{chardon1991} A.~Chardon, I.~Cretois, C.~Hourseau.
\emph{Skin colour typology and suntanning pathways}. International Journal of
Cosmetic Science, 13(4):191--208, 1991. (ITA)

\bibitem{nowara2020} E.~M.~Nowara, D.~McDuff, A.~Veeraraghavan.
\emph{A meta-analysis of the impact of skin tone and gender on non-contact
photoplethysmography measurements}. CVPR Workshops, 2020.

\bibitem{dasari2021} A.~Dasari, S.~K.~A.~Prakash, L.~A.~Jeni, C.~S.~Tucker.
\emph{Evaluation of biases in remote photoplethysmography methods}. npj Digital
Medicine, 4:91, 2021.

\bibitem{toye2023} P.~Toye. \emph{VitalVideos: A dataset of face videos with PPG
and blood pressure ground truth}. 2023. \url{https://vitalvideos.org}

\end{thebibliography}

\end{document}
